\documentclass[11pt]{article}
\usepackage[margin=1in]{geometry}
\usepackage{amsmath,amssymb,amsthm,mathtools}
\usepackage{enumitem}
\usepackage{hyperref}
\usepackage[T1]{fontenc}

\newtheorem{theorem}{Theorem}[section]
\newtheorem{proposition}[theorem]{Proposition}
\newtheorem{corollary}[theorem]{Corollary}
\newtheorem{conjecture}[theorem]{Conjecture}
\theoremstyle{definition}
\newtheorem{definition}[theorem]{Definition}
\newtheorem{remark}[theorem]{Remark}

\title{Resource-Localized Obstruction Profiles for Circuit Lower Bounds:\\
Formal Calibration Theorems and Exact Remaining Gaps}
\author{ChatGPT}
\date{September 2026}

\begin{document}
\maketitle

\begin{abstract}
We develop a resource-sensitive refinement of Verifier-Invariant Obstruction
Theory for the study of circuit lower bounds.  The objective is not to attach
homology to one arbitrary verifier, but to construct quantitative obstructions
only after quotienting computational presentations by explicitly controlled
simulations.  We prove and Lean-formalize five calibration results.  First,
same-language equivalence together with constant-factor equivalence of cost
profiles defines a quotient of costed verifier presentations.  Second, a pure
semantic quotient provably forgets asymptotic cost.  Third, bidirectional
certificate simulations with constant-factor overhead form an equivalence
relation, and ignored certificate padding is trivial in the corresponding
quotient.  Fourth, raw trace-cycle structure changes under constant-factor
resource equivalence.  Fifth, a potential whose increase per construction
step is bounded yields a quantitative construction-size lower bound, and a
family-level superpolynomial potential-growth condition rules out polynomially
bounded derivations.

We also prove that the canonical pointwise-disagreement obstruction is
logically equivalent to the desired lower-bound statement and therefore
provides no independent leverage.  These results motivate Resource-Localized
Obstruction Profiles (RLOPs).  The decisive missing theorem remains a
non-tautological RLOP for unrestricted Boolean circuits with superpolynomial
growth on an explicit NP-complete family.  No resolution of P versus NP is
claimed.
\end{abstract}

\section{Introduction}

A circuit-lower-bound route to \(P\ne NP\) may target the stronger statement
\[
  \mathrm{SAT}\notin P/poly.
\]
To prove it by an obstruction \(\Omega\), one needs substantially more than a
cycle, cohomology class, or geometric feature of one machine presentation.
One needs a quantitative theorem of the form
\[
  C\text{ computes }f
  \quad\Longrightarrow\quad
  \Omega(f,C)\le B(|C|),
\]
and a separate explicit-target theorem showing that every polynomial-size
candidate violates the required bound.

Raw computation graphs are not canonical.  Certificate padding, state
recoding, bounded stuttering, dummy branching, and other constant-overhead
changes can alter their combinatorial and topological structure while
preserving both the language and its asymptotic resource class.  Thus a
candidate invariant must first descend through a resource-controlled
localization.

The present paper gives formal calibration theorems for this principle.  The
results are machine-checked in the companion Lean development.  They narrow
the design space for a valid obstruction but do not establish a lower bound
for unrestricted SAT circuits.

\section{Costed verifier presentations}

\begin{definition}[Cost profile]
A cost profile is a function
\[
  T:\mathbb N\to\mathbb N.
\]
\end{definition}

\begin{definition}[Constant-factor equivalence]
For cost profiles \(T,U\), write \(T\asymp_c U\) when there is a constant
\(K\in\mathbb N\) such that for every \(n\),
\[
  T(n)\le K U(n)
  \quad\text{and}\quad
  U(n)\le K T(n).
\]
\end{definition}

\begin{proposition}
Constant-factor equivalence is an equivalence relation.
\end{proposition}

\begin{proof}
Reflexivity uses \(K=1\).  Symmetry exchanges the two inequalities.  For
transitivity, suppose \(T\asymp_c U\) with factor \(K\) and
\(U\asymp_c V\) with factor \(L\).  Then
\[
  T(n)\le K U(n)\le KL V(n)
\]
and, in the reverse direction,
\[
  V(n)\le L U(n)\le LK T(n)=KL T(n).
\]
\end{proof}

\begin{definition}[Resource presentation]
A resource presentation over an input type \(X\) is a triple
\[
  \mathcal V=(W,V,T),
\]
where \(W\) is a certificate type,
\(V:X\times W\to\{0,1\}\) is a verifier, and \(T\) is a cost profile.
Its language is
\[
  L(\mathcal V)=\{x:\exists w,\ V(x,w)=1\}.
\]
\end{definition}

\begin{definition}[First resource equivalence]
Two resource presentations are first-resource equivalent when they recognize
the same language and their cost profiles are constant-factor equivalent.
\end{definition}

\begin{theorem}
First-resource equivalence is an equivalence relation and therefore defines a
quotient of resource presentations.
\end{theorem}

\begin{proof}
Same-language equivalence and constant-factor equivalence are both equivalence
relations.  Their conjunction is again an equivalence relation.
\end{proof}

This quotient is deliberately a calibration object.  Its cost annotation is
supplied externally rather than derived from a formal operational semantics.
Nevertheless, it already distinguishes semantic equivalence from resource
equivalence.

\section{The semantic quotient forgets cost}

\begin{theorem}[Semantic cost-loss theorem]
There exist two resource presentations whose underlying verifier
presentations have the same semantic class but whose cost profiles are not
constant-factor equivalent.
\end{theorem}

\begin{proof}
Use the same verifier in both presentations: on the sole input and sole
certificate, always return true.  Assign the first presentation the cost
profile
\[
  T_1(n)=1
\]
and the second the profile
\[
  T_2(n)=n+1.
\]
After cost is forgotten, the verifier presentations are literally identical.
Suppose, however, that the profiles were constant-factor equivalent with
factor \(K\).  The reverse comparison would imply
\[
  n+1\le K
\]
for every \(n\).  Setting \(n=K\) gives \(K+1\le K\), a contradiction.
\end{proof}

\begin{corollary}
An obstruction defined only on the pure semantic quotient cannot, without
additional data, recover asymptotic running time or circuit size.
\end{corollary}

The conclusion does not say that semantic invariants are useless.  It says
that a complexity obstruction must be enriched by resource data or by a
resource-controlled morphism theory.

\section{Explicit resource simulations}

\begin{definition}[Certificate simulation]
Let \(\mathcal V=(W,V,T)\) and \(\mathcal U=(Z,U,S)\).  A certificate
simulation from \(\mathcal V\) to \(\mathcal U\) is a map
\[
  \phi:W\to Z
\]
such that for every input \(x\) and certificate \(w\),
\[
  U(x,\phi(w))=V(x,w).
\]
\end{definition}

Such a simulation gives the language inclusion
\[
  L(\mathcal V)\subseteq L(\mathcal U).
\]
Indeed, an accepting certificate \(w\) maps to an accepting certificate
\(\phi(w)\).

\begin{definition}[Resource simulation]
A resource simulation additionally requires a constant \(K\) for which
\[
  S(n)\le K T(n)
\]
for all \(n\).  A resource bisimulation consists of resource simulations in
both directions.
\end{definition}

\begin{theorem}
Resource bisimulation is an equivalence relation.
\end{theorem}

\begin{proof}
Identity certificate maps and the factor one give reflexivity.  Reversing the
two simulations gives symmetry.  Certificate maps compose, verifier-answer
preservation is transitive, and cost constants multiply, giving
transitivity.
\end{proof}

\begin{theorem}[Padding is trivial in the simulation quotient]
Let \(D\) be nonempty and define the padded verifier by
\[
  V_D(x,(d,w))=V(x,w).
\]
If padding leaves the supplied cost profile unchanged, then \(\mathcal V\)
and \(\mathcal V_D\) are resource-bisimilar.
\end{theorem}

\begin{proof}
Choose one element \(d_0\in D\).  The forward certificate map is
\[
  w\longmapsto(d_0,w),
\]
and the backward map is
\[
  (d,w)\longmapsto w.
\]
Both maps preserve the verifier answer definitionally, and both cost
comparisons use factor one.
\end{proof}

\begin{corollary}
Every invariant defined on the resource-simulation quotient ignores nonempty
certificate padding.
\end{corollary}

\section{Trace topology is not resource invariant}

Consider a direct computation graph
\[
  s\longrightarrow t
\]
and a diamond graph
\[
  s\longrightarrow \ell\longrightarrow t,
  \qquad
  s\longrightarrow r\longrightarrow t.
\]
Both graphs are acyclic and have strictly increasing rank functions.  Under
the fixed-start/fixed-accept abstract trace semantics, both recognize the
universal language on a one-element input type.

\begin{theorem}[Reachable trace universality]
For the abstract language
\[
  L_{G,s,t}(x)\iff s\leadsto_G t,
\]
if \(s\leadsto_G t\), then \(L_{G,s,t}\) is the universal language.
Consequently, any two reachable trace presentations recognize the same
language.
\end{theorem}

\begin{proof}
The right-hand side of the defining equivalence is independent of \(x\).  If
the fixed reachability fact holds, it witnesses membership for every input.
\end{proof}

Equip the direct graph with constant cost one and the diamond graph with
constant cost two.  The two profiles are constant-factor equivalent.
Nevertheless, after forgetting edge directions the diamond has the induced
four-cycle
\[
  s-\ell-t-r-s,
\]
while the two-vertex direct graph has no four distinct vertices.

\begin{theorem}[Resource-sensitive trace no-go theorem]
Induced-four-cycle structure is not invariant under same-language,
constant-factor resource equivalence.
\end{theorem}

\begin{proof}
The direct and diamond presentations have the same language and equivalent
constant profiles, but their induced-four-cycle predicates have opposite
truth values.
\end{proof}

This counterexample uses neither unreachable states nor superpolynomial
padding.  It follows that raw graph cycles cannot serve as resource-sensitive
language obstructions before quotienting constant-overhead branch insertion
and merging.

\section{Quantitative potential bridges}

\begin{definition}[Step derivation]
Let \(S\to S'\) denote one legal construction step.  A derivation of length
\(m\) is a sequence of exactly \(m\) such steps.
\end{definition}

\begin{theorem}[Potential accumulation]
Let \(\Phi\) be a natural-valued potential.  Suppose every legal step
satisfies
\[
  \Phi(S')\le \Phi(S)+\delta.
\]
Then every length-\(m\) derivation from \(S_0\) to \(S_m\) satisfies
\[
  \Phi(S_m)\le \Phi(S_0)+m\delta.
\]
\end{theorem}

\begin{proof}
Induct on the derivation.  The zero-step case is equality.  For the last
step, combine the induction hypothesis with the one-step bound:
\[
  \Phi(S_{m+1})
  \le \Phi(S_m)+\delta
  \le \Phi(S_0)+m\delta+\delta
  =\Phi(S_0)+(m+1)\delta.
\]
\end{proof}

\begin{corollary}
If
\[
  \Phi(S_0)+m\delta<\Phi(S_{\mathrm{target}}),
\]
then no length-\(m\) derivation reaches the target.
\end{corollary}

For a family indexed by input length, let \(m(n)\) be the construction size,
\(\delta(n)\) the one-step growth bound, and \(S_0(n),S_f(n)\) the initial
and target states.

\begin{theorem}[Family-level superpolynomial bridge]
Assume that for every coefficient \(a\) and degree \(k\), some \(n\) satisfies
\[
 \Phi_n(S_0(n))+
 \bigl(a(n+1)^k\bigr)\delta(n)
 <
 \Phi_n(S_f(n)).
\]
If every target has a derivation of size \(m(n)\), then \(m\) is not
polynomially bounded.
\end{theorem}

\begin{proof}
Suppose \(m(n)\le a(n+1)^k\) for fixed \(a,k\).  The potential accumulation
theorem gives
\[
 \Phi_n(S_f(n))
 \le \Phi_n(S_0(n))+m(n)\delta(n)
 \le \Phi_n(S_0(n))+a(n+1)^k\delta(n),
\]
contradicting the assumed target-growth inequality for the corresponding
\(n\).
\end{proof}

For circuit lower bounds, one would instantiate a state with a library of
available gate outputs and a construction step with the addition of one legal
gate.  The theorem is a correct circuit-soundness schema.  Its force depends
entirely on finding a non-tautological potential with a strong local bound and
superpolynomial target value.

\section{The disagreement diagnostic}

\begin{definition}
For Boolean functions \(C,f:X\to\{0,1\}\), define
\[
  D(C,f)\iff \exists x,\ C(x)\ne f(x).
\]
\end{definition}

\begin{theorem}
For every candidate class \(\mathcal C\),
\[
  \forall C\in\mathcal C,\ D(C,f)
  \quad\Longleftrightarrow\quad
  \nexists C\in\mathcal C\text{ computing }f.
\]
\end{theorem}

\begin{proof}
If every candidate disagrees, no candidate is exact.  Conversely, if a
candidate had no disagreement, pointwise equality would make it an exact
computation, contradicting the absence of an exact candidate.
\end{proof}

\begin{remark}
The theorem is useful as an audit.  Choosing disagreement itself as
\(\Omega\) makes ``target nonvanishing'' logically identical to the desired
lower bound.  It is not an independently analyzable obstruction.
\end{remark}

\section{Resource-Localized Obstruction Profiles}

\begin{definition}[RLOP]
A Resource-Localized Obstruction Profile for a target family \(f_n\) should
include:
\begin{enumerate}[leftmargin=*]
\item a quotient or localization by explicitly bounded computational
simulations;
\item a gate-construction category, rewriting system, or higher groupoid;
\item a presentation-invariant potential or index \(\Phi_{f_n}\);
\item a theorem bounding the change caused by one unrestricted circuit gate;
\item a superpolynomial lower bound for \(\Phi_{f_n}\) on the explicit target;
\item a uniformity and model-transfer theorem; and
\item an audit against relativization, natural proofs, and algebrization.
\end{enumerate}
\end{definition}

The calibration theorems in this paper establish the logical form of items
one and four, and show why a semantic quotient or raw trace invariant is
insufficient.  They do not produce the potential required by items three and
five.

\section{Barrier audit}

The quotient and potential-accumulation theorems are deliberately generic.
They relativize and remain valid under algebraic extensions.  Therefore they
cannot be the decisive nonrelativizing or nonalgebrizing step in a proof of
\(P\ne NP\).

A candidate-relative obstruction may avoid the largeness premise of the
Razborov--Rudich framework because it depends jointly on the target and the
candidate circuit.  Yet an efficient refuter for every unrestricted small
circuit would itself be a major constructive lower bound.  Thus the phrase
``candidate relative'' does not remove the central difficulty.

Finally, a theorem for monotone, bounded-depth, threshold, formula,
noncommutative arithmetic, or proof-complexity models must not be promoted to
unrestricted Boolean circuits without an explicit simulation theorem and a
proved overhead bound.

\section{Exact remaining target}

A complete RLOP separation must produce an explicit NP-complete family,
preferably \(\mathrm{SAT}_n\), and prove all of the following:
\begin{enumerate}[leftmargin=*]
\item every polynomial-size unrestricted circuit induces a construction with
polynomially many legal steps;
\item the RLOP potential changes by at most a controlled amount per gate;
\item the initial potential is controlled;
\item the SAT target potential exceeds every polynomial construction budget;
\item the construction is uniform enough to identify an NP language;
\item the decisive target-growth proof escapes the classical barriers.
\end{enumerate}
The family-level potential theorem would then yield
\(\mathrm{SAT}\notin P/poly\), and the standard inclusion
\(P\subseteq P/poly\) would imply \(P\ne NP\).

At present, the fourth and sixth items are unresolved.  They are not hidden
inside the terminology of RLOPs.

\section{Lean formalization}

The companion modules formalize all proved calibration results without
`Sorry', `admit', or user-declared axioms:
\begin{itemize}[leftmargin=*]
\item `ResourceSensitivePresentation.lean';
\item `SemanticCostLoss.lean';
\item `ResourceSimulation.lean';
\item `TracePresentationUniversality.lean';
\item `ResourceSensitiveTraceNoGo.lean';
\item `QuantitativePotential.lean';
\item `AsymptoticPotential.lean'; and
\item `DisagreementObstruction.lean'.
\end{itemize}

\section{Conclusion}

Resource sensitivity is not a cosmetic addition to topological circuit
lower bounds.  Pure semantic quotients erase cost, while raw trace topology
changes under constant-overhead presentation transformations.  A valid
obstruction must therefore descend through an explicit simulation quotient
and simultaneously obey a quantitative gate-growth theorem.

The RLOP framework isolates these requirements and supplies a formal
family-level bridge from target growth to superpolynomial construction size.
The unresolved scientific core is now explicit: construct a non-tautological
RLOP over unrestricted Boolean gates and prove superpolynomial growth for an
explicit NP-complete family in a way that escapes the known barriers.

\begin{thebibliography}{99}
\bibitem{BGS}
T. Baker, J. Gill, and R. Solovay,
Relativizations of the \(P=?NP\) question,
\emph{SIAM Journal on Computing} 4 (1975).

\bibitem{RR}
A. Razborov and S. Rudich,
Natural proofs,
\emph{Journal of Computer and System Sciences} 55 (1997).

\bibitem{AW}
S. Aaronson and A. Wigderson,
Algebrization: A new barrier in complexity theory,
\emph{ACM Transactions on Computation Theory} 1 (2009).

\bibitem{Williams}
R. Williams,
Nonuniform ACC circuit lower bounds,
\emph{Journal of the ACM} 61 (2014).

\bibitem{LiYang}
J. Li and T. Yang,
A \(3.1n-o(n)\) circuit lower bound for explicit functions,
STOC, 2022.

\bibitem{CDJ}
M. Carmosino, N. Dang, and T. Jackman,
Convergent gate elimination and constructive circuit lower bounds,
arXiv:2602.17942, 2026.

\bibitem{CJSSW}
L. Chen, C. Jin, R. Santhanam, and R. Williams,
Constructive separations and their consequences,
ECCC TR21-159, 2021.
\end{thebibliography}

\end{document}
